\section*{Capítulo \ref{ch:b3:ode} --- Equações diferenciais
ordinárias}

\begin{solution}{exo:b3:ode:1}
(a) Separando as variáveis: $x(t) = \frac1{1 - t}$ em
$\intoo{-\infty}1$: explosão em $T_+ = 1$, com $x(t) \to
+\infty$ --- exatamente o~\cref{thm:b3:ode:maximal}(2).
(b) $x(t) = \tan t$ em $\intoo{-\pi/2}{\pi/2}$: explosão nos
dois extremos.
(c) $x(t) = \frac{1}{1 + \eu^{-t}}$, global: a solução
permanece em $\intoo01$, um conjunto limitado, de modo que o gráfico não pode
escapar de todo \omterm{def:b3:topology:compact}{compacto} de $\R\times\R$ em tempo finito --- $T_\pm
= \pm\infty$. Note que (a) e (b) não contradizem
o~\cref{cor:b3:ode:globallinear}: $x^2$ e $1 + x^2$ têm
crescimento superlinear.
\end{solution}

\begin{solution}{exo:b3:ode:2}
(a) Essa é a estimativa do~\cref{cor:b3:ode:uniqueness} com $t_0
= 0$. Otimalidade: para $x' = Lx$ (globalmente $L$-lipschitziana), duas
soluções diferem exatamente de $(x_0 - y_0)\eu^{Lt}$.
(b) A estimativa se lê: a aplicação de \omterm{ex:b3:ode:planeclassification}{fluxo} no tempo $t$ é
$\eu^{L\abs t}$-lipschitziana na condição inicial ---
\omterm{def:b3:topology:continuity}{continuidade}, uniformemente para $t$ em \omterm{def:b3:topology:compact}{compactos}; em conjuntos limitados de
$f$ não globalmente lipschitzianas, repita o mesmo num tubo \omterm{def:b3:topology:compact}{compacto}
em torno das trajetórias com a constante local.
\end{solution}

\begin{solution}{exo:b3:ode:3}
(a) $\abs{\sin(tx)} \leq 1$: limitada, isto é, crescimento linear com
$\alpha = 0$, $\beta = 1$: \cref{cor:b3:ode:globallinear} em
$U = \R\times\R$.
(b) $\abs{tx/(1 + x^2)} \leq \abs t\cdot\frac12$: de novo
sublinear (de fato limitada em \omterm{def:b3:topology:compact}{compactos} de tempo): global.
(c) $X = (x, x')$: $X' = \bigl(\begin{smallmatrix}0 & 1\\ -q(t)
& 0\end{smallmatrix}\bigr)X$: linear com coeficientes
\omterm{def:b3:topology:continuity}{contínuos}: global (o quadro do~\cref{thm:b3:ode:linearstructure}).
(d) Global; Grönwall como no~\cref{cor:b3:ode:globallinear}:
$\norm{x(t)} \leq \norm{x(t_0)}\,\eu^{M\abs{t - t_0}}$ com
$M = \sup\vertiii{A}$.
\end{solution}

\begin{solution}{exo:b3:ode:4}
(a) $A^2 = -I$ para a primeira: $\eu^{tA} = \cos t\,I + \sin
t\,A = \bigl(\begin{smallmatrix}\cos t & -\sin t\\ \sin t &
\cos t\end{smallmatrix}\bigr)$. Bloco de Jordan: $\lambda I$ e
$N$ comutam: $\eu^{tA} =
\eu^{\lambda t}\bigl(\begin{smallmatrix}1 & t\\ 0 &
1\end{smallmatrix}\bigr)$. Terceira: $A^2 = I$: $\eu^{tA} =
\cosh t\,I + \sinh t\,A$.

(b) Sistema $X' = \bigl(\begin{smallmatrix}0&1\\-1&0
\end{smallmatrix}\bigr)X + \bigl(\begin{smallmatrix}0\\
\cos\omega t\end{smallmatrix}\bigr)$; Duhamel com o
\omterm{thm:b3:ode:linearstructure}{resolvente} da rotação dá as soluções particulares: para
$\omega \neq 1$, $x_p(t) = \frac{\cos(\omega t)}{1 -
\omega^2}$ (verifique diretamente); para $\omega = 1$, a integral
$\int_0^t\sin(t - s)\cos s\,\dd s = \frac t2\sin t$ produz
o crescimento \emph{secular} $x_p = \frac t2\sin t$: ressonância
--- a forçante bombeia energia na frequência natural e a
amplitude cresce linearmente.
\end{solution}

\begin{solution}{exo:b3:ode:5}
(a) $w' = x_1x_2'' - x_1''x_2 = x_1(-px_2' - qx_2) - (-px_1'
- qx_1)x_2 = -p\,w$: $w(t) = w(t_0)\eu^{-\int_{t_0}^tp}$,
nunca nulo nem identicamente nulo. Se $w \neq 0$, os vetores
$(x_i, x_i')(t_0)$ são independentes em $\R^2$ e, como o
espaço de soluções tem dimensão $2$
(\cref{thm:b3:ode:linearstructure} para o sistema), $(x_1,
x_2)$ é uma base.

(b) Com $u = \int\frac{\eu^{-\int p}}{x_1^2}$: $x_2 = x_1u$,
$x_2' = x_1'u + \frac{\eu^{-\int p}}{x_1}$, e
\[
x_2'' + px_2' + qx_2
= u\,(x_1'' + px_1' + qx_1) +
\Bigl(-\,p\frac{\eu^{-\int p}}{x_1} +
p\frac{\eu^{-\int p}}{x_1}\Bigr) = 0
\]
(os termos cruzados se cancelam exatamente; expanda com cuidado). Para
$t^2x'' - 2x = 0$, isto é, $x'' - \frac2{t^2}x = 0$ ($p = 0$)
com $x_1 = t^2$: $u = \int t^{-4} = -\frac1{3t^3}$, de modo que $x_2
= -\frac1{3t}$: solução geral $at^2 + \frac bt$.
\end{solution}

\begin{solution}{exo:b3:ode:6}
Equilíbrios $0, 1$; $f'(x) = 1 - 2x$: $f'(0) = 1 > 0$
(instável --- as soluções próximas se afastam, como a forma
explícita mostra), $f'(1) = -1 < 0$: assintoticamente \omterm{def:b3:ode:stability}{estável}
(\cref{thm:b3:ode:linearization} em dimensão $1$). Para
$x(0) \in \intoo01$: $f > 0$ ali, de modo que, enquanto a
solução permanecer em $\intoo01$, ela cresce; ela nunca pode atingir
$0$ nem $1$ (unicidade: essas são trajetórias), de modo que ela permanece,
é limitada --- logo global --- e cresce até um limite $L
\in \intoc{x(0)}1$. Se $f(L) \neq 0$, então $x' \geq c > 0$
perto do limite, forçando $x$ a passar de $L$: logo $f(L) = 0$, $L = 1$;
simetricamente, $x \to 0$ em $-\infty$. Explicitamente, $x(t) =
\frac1{1 + C\eu^{-t}}$ confirma tudo. Em geral: se uma
solução escalar autônoma tivesse $x'(t_0) = 0$, então $x(t_0)$
é um equilíbrio e a unicidade torna $x$ constante;
caso contrário, $f(x(t))$ mantém sinal fixo (ela nunca se anula,
e $t \mapsto f(x(t))$ é \omterm{def:b3:topology:continuity}{contínua}): $x$ é estritamente
monótona --- de modo que uma solução periódica não constante é
impossível.
\end{solution}

\begin{solution}{exo:b3:ode:7}
(a) $\frac{\dd}{\dd t}H(x,y) = H_xx' + H_yy' = H_xH_y +
H_y(-H_x) = 0$.
(b) Os conjuntos de nível de $H = \frac{y^2}2 + \frac{x^4}4$ são
\omterm{def:b3:topology:compact}{compactos} ($H$ coerciva), de modo que as soluções ficam presas em \omterm{def:b3:topology:compact}{compactos}:
globais e limitadas (\cref{thm:b3:ode:maximal}). Estabilidade de
$(0,0)$: $V = H$ é definida positiva ($H = 0$ só na
origem) com $\dot V = 0$: \cref{thm:b3:ode:lyapunov}. A
linearização $x' = y$, $y' = 0$ tem a matriz nilpotente
não diagonalizável de autovalor $0$:
o~\cref{thm:b3:ode:linearization} nada diz (sua hipótese
$\operatorname{Re}\lambda < 0$ falha) e, de fato, o
sistema linearizado é instável ($y_0 \ne 0$ deriva), ao passo que o
não linear é \omterm{def:b3:ode:stability}{estável}: a linearização num equilíbrio não
hiperbólico nada demonstra.
\end{solution}

\begin{solution}{exo:b3:ode:8}
(a) $\dot V = y\,y' + \sin x\cdot x' = y(-\sin x - cy) +
y\sin x = -cy^2 \leq 0$, e $V = \frac{y^2}2 + (1 - \cos x)$
é definida positiva em $\{\abs x < 2\pi\}$ em torno da
origem: \omterm{def:b3:ode:stability}{estável} (\cref{thm:b3:ode:lyapunov}).
(b) A matriz linearizada em $(0,0)$ é
$\bigl(\begin{smallmatrix}0 & 1\\ -1 &
-c\end{smallmatrix}\bigr)$, de polinômio característico
$\lambda^2 + c\lambda + 1$: raízes $\frac{-c \pm \sqrt{c^2 -
4}}2$ --- ambas reais negativas se $c \geq 2$, complexas de parte
real $-\frac c2 < 0$ se $0 < c < 2$. Em todos os casos
$\operatorname{Re}\lambda < 0$:
o~\cref{thm:b3:ode:linearization} dá estabilidade assintótica
(apesar do $\dot V$ degenerado).
(c) Em $(\pi, 0)$: $\sin(\pi + u) = -\sin u$, linearização
$\bigl(\begin{smallmatrix}0&1\\1&-c\end{smallmatrix}\bigr)$,
característico $\lambda^2 + c\lambda - 1$: raízes de sinais
opostos ($\lambda_+\lambda_- = -1$). Ao longo do autovetor
instável, o sistema linear tem a solução explicitamente fugidia
$\eu^{\lambda_+t}v_+$ com $\lambda_+ > 0$: o
pêndulo invertido é instável para todo amortecimento.
\end{solution}

\begin{solution}{exo:b3:ode:9}
Para $\alpha \in \intoo01$: além de $x \equiv 0$, cada
\[
x_c(t) = \begin{cases}0 & t \leq c,\\
\bigl((1-\alpha)(t - c)\bigr)^{1/(1-\alpha)} & t \geq c,
\end{cases}
\]
é $\mathcal C^1$ e resolve a equação (o expoente
$\frac1{1-\alpha} > 1$ faz a derivada se anular em $c$): um
\omterm{def:b3:topology:continuity}{contínuo} de soluções por $(0,0)$. Para $\alpha = 1$: $x
\mapsto \abs x$ é globalmente $1$-lipschitziana
($\abs{\abs a - \abs b} \leq \abs{a - b}$), de modo que
o~\cref{thm:b3:ode:cauchylipschitz} se aplica e a única
solução por $0$ é $x \equiv 0$. A fronteira é exatamente
a condição de Lipschitz local em $0$: $\abs x^\alpha$ tem
quocientes de diferenças ilimitados ali para $\alpha < 1$.
\end{solution}

\begin{solution}{exo:b3:ode:10}
(a) Suponha $\norm{x(t_2)} > R$ para algum $t_2 > 0$ com
$\norm{x(0)} \leq R$, e seja $t_1 = \sup\{t \leq t_2 :
\norm{x(t)} \leq R\}$: então $\norm{x(t_1)} = R$ e
$\norm{x(t)} > R$ em $\intoc{t_1}{t_2}$. Nesse intervalo,
$g(t) = \norm{x(t)}^2$ tem $g'(t) = 2\langle x(t),
F(x(t))\rangle \leq 0$ (a hipótese se aplica: $\norm{x(t)}
\geq R$), de modo que $g(t_2) \leq g(t_1) = R^2$: contradição. A
bola é positivamente invariante.
(b) Uma solução presa na bola compacta não pode ter
$T_+ < \infty$ (\cref{thm:b3:ode:maximal}(2)): global
para a frente. Sistema gradiente: $\frac{\dd}{\dd t}G(x(t)) =
\langle\nabla G, -\nabla G\rangle = -\norm{\nabla G(x(t))}^2
\leq 0$: $G$ decresce, de modo que a solução permanece em $\{G \leq
G(x_0)\}$, que é limitado (coercividade: fora de uma bola
grande, $G > G(x_0)$) e fechado: \omterm{def:b3:topology:compact}{compacto}. A fuga é
impossível: toda solução de um sistema gradiente coercivo é
global para a frente, deslizando ladeira abaixo para sempre.
\end{solution}

\begin{solution}{exo:b3:ode:11}
(a) Lema de comparação: seja $w = x - y$ no intervalo
comum; $w(0) \geq 0$ e $w' = x' - y' \geq F(x) - F(y) =
c(t)w$ com $c(t) = \frac{F(x) - F(y)}{x - y}$ limitada em
intervalos de tempo \omterm{def:b3:topology:compact}{compactos} ($F$ \omterm{def:b3:ode:lipschitz}{localmente lipschitziana}); então
$(w\eu^{-\int c})' \geq 0$, de modo que $w \geq 0$ em toda parte. Com
$F(x) = x^2$: $y(t) = \frac1{1 - t}$ explode em $1$, e $x
\geq y$ enquanto $x$ viver; se $T_+ > 1$, então $x$ seria
finita em $t = 1$ enquanto domina $y \to \infty$:
absurdo. $T_+ \leq 1$.

(b) A comparação invertida (mesmo lema, papéis trocados):
em $\intcc01\cap\intco0{T_+}$, $t^2 \leq 1$ dá $x' \leq
x^2 + 1$, ao passo que $z(t) = \tan(t + \frac\pi4)$ satisfaz $z' =
z^2 + 1$, $z(0) = 1 = x(0)$: logo $x \leq z$ enquanto
ambas estiverem definidas. Como $z$ é finita em
$\intco0{\frac\pi4}$, $x$ não pode explodir antes de
$\frac\pi4$: $T_+ \geq \frac\pi4$.

(c) Juntas: $\frac\pi4 \leq T_+ \leq 1$ (numericamente $T_+
\approx 0.96$). Moral: para $x' = F(t, x)$ com $F$
superlinear em $x$, as soluções explodem em tempo finito
sempre que $\int^\infty\frac{\dd s}{F(s)} < \infty$ (a
solução de comparação atinge o infinito nesse tempo finito);
o crescimento linear, em que a integral diverge, força a existência
global (\cref{exo:b3:ode:3}). É a mesma integral de Osgood
do~\cref{pb:b3:complete:1}, governando agora a
fuga para o infinito em vez da fuga a partir do zero.
\end{solution}

\begin{solution}{exo:b3:ode:12}
(a) $W' = uv'' - u''v = u(-q_2v) - (-q_1u)v = (q_1 -
q_2)\,uv$.

(b) Sejam $a < b$ zeros consecutivos de $u$; normalize $u >
0$ em $\intoo ab$ (de modo que $u'(a) > 0$, $u'(b) < 0$ --- não nulos
por unicidade, já que $u(a) = u'(a) = 0$ forçaria $u
\equiv 0$). Suponha que $v$ não tenha zero em $\intoo ab$;
normalize $v > 0$ ali (logo $v(a), v(b) \geq 0$ por
\omterm{def:b3:topology:continuity}{continuidade}). Integre (a):
\[
W(b) - W(a) = \int_a^b(q_1 - q_2)\,uv \;\leq\; 0 .
\]
Mas $W(a) = u(a)v'(a) - u'(a)v(a) = -u'(a)v(a) \leq 0$ e
$W(b) = -u'(b)v(b) \geq 0$: logo $W(b) - W(a) \geq 0$.
Igualdade em toda a cadeia: $\int(q_1 - q_2)uv = 0$ com $uv > 0$
no intervalo aberto força $q_1 = q_2$ ali; e $W(a) =
W(b) = 0$ força $v(a) = v(b) = 0$; então $W \equiv 0$ em
$\intcc ab$ (sua derivada se anula), isto é,
$(v/u)' = -W/u^2 = 0$ em $\intoo ab$: $v \propto u$.

(c) Tome $q_1 = m$ e $u = \sin(\sqrt m(t - t_0))$, cujos
zeros consecutivos distam $\pi/\sqrt m$, e $q_2 = q \geq
m$: por (b), toda solução $v$ de $v'' + qv = 0$ se anula em
cada intervalo aberto de comprimento $\pi/\sqrt m$ (na
alternativa degenerada $q \equiv m$, $v \propto u$ também se anula). Se, em vez disso, $q \leq 0$: aplique (b) com $q_1 = q$, $u =
v$ e $q_2 = 0$ com a solução sem zeros $\mathbf 1$ de
$v'' = 0$. Se $v$ tivesse dois zeros consecutivos, (b) forçaria
ou um zero de $\mathbf 1$ entre eles, ou o caso degenerado
$\mathbf 1 \propto v$ --- ambos absurdos: $v$ se anula no
máximo uma vez. Testes: para $q = 1$, $\sin t$ se anula a cada $\pi
= \pi/\sqrt1$, como previsto; para $q = -1$, $\sinh t$
se anula exatamente uma vez e $\eu^t$ nunca --- no máximo um
zero, como previsto.
\end{solution}

\begin{solution}{pb:b3:ode:1}
\textbf{1.} $\dot E = yy' + \sin x\cdot x' = -y\sin x +
y\sin x = 0$: a energia é uma integral primeira. Numa
\omterm{thm:b3:ode:maximal}{solução maximal}, $y^2 = 2(E_0 + \cos x) \leq 2(E_0 + 1)$: $y$ é
limitada; então $\abs{x(t)} \leq \abs{x(0)} +
t\sup\abs y$ cresce no máximo linearmente: em qualquer intervalo
de tempo finito a trajetória permanece num \omterm{def:b3:topology:compact}{compacto} de $\R^2$, de modo que
o~\cref{thm:b3:ode:maximal}(2) força $T_\pm = \pm\infty$.
Equilíbrios $(k\pi, 0)$; linearização
$\bigl(\begin{smallmatrix}0&1\\
\mp1&0\end{smallmatrix}\bigr)$ com $-\cos(k\pi) = \mp1$:
autovalores $\pm\iu$ para $k$ par (tipo \omterm{ex:b3:groups:actions}{centro}, inconclusivo
por si só) e $\pm1$ para $k$ ímpar (sela).

\textbf{2.} $E$ constante ao longo das soluções confina cada
trajetória a um conjunto de nível $\{y^2 = 2(E_0 + \cos x)\}$. Para
$E_0 = -1$: apenas os pontos $(2k\pi, 0)$. Para $-1 < E_0 <
1$: escrevendo $E_0 = -\cos a$, o conjunto é uma união disjunta de
curvas fechadas $y = \pm\sqrt{2(\cos x - \cos a)}$ sobre $x \in
[2k\pi - a, 2k\pi + a]$, uma em torno de cada equilíbrio \omterm{def:b3:ode:stability}{estável}.
Para $E_0 = 1$: as curvas $y = \pm2\cos\frac x2$ que ligam
selas consecutivas --- a separatriz --- junto com as
próprias selas. Para $E_0 > 1$: dois gráficos $y =
\pm\sqrt{2(E_0 + \cos x)}$, definidos para todo $x$, nunca
tocando $y = 0$.

\textbf{3.} $V = E + 1 = \frac{y^2}2 + (1 - \cos x)$
se anula em $(0,0)$, é positiva numa \omterm{def:b3:topology:topology}{vizinhança} perfurada
($\abs x < 2\pi$), e $\dot V = 0 \leq 0$:
o~\cref{thm:b3:ode:lyapunov} dá a estabilidade. Não assintótica:
a solução por $(a, 0)$ ($0 < a$ pequeno) permanece na
curva de nível $E = -\cos a$, cuja distância à origem é
positiva (a curva encontra o eixo $x$ apenas em $\pm a$):
$\varphi_t(a, 0) \not\to (0,0)$.

\textbf{4.} Na curva de nível $C_a$: sem equilíbrios ($y = 0$
força $x = \pm a$ com $\sin(\pm a) \neq 0$ para $0 < a <
\pi$), de modo que a velocidade $\norm{(y, -\sin x)}$ tem mínimo positivo
$m$ no \omterm{def:b3:topology:compact}{compacto} $C_a$. Siga a solução a partir de
$(a, 0)$: no semiplano inferior $x' = y < 0$, de modo que $x$
decresce de $a$ a $-a$ em tempo finito (as integrais de
quarto/meio período convergem: a análise da questão 5), atingindo
$(-a, 0)$; pela simetria $(x, y) \mapsto (x, -y)$, $t
\mapsto -t$ da equação, a metade superior é percorrida
de volta no mesmo tempo $\frac T2$: a solução retorna a
$(a, 0)$ no tempo $T$. A unicidade
(\cref{cor:b3:ode:uniqueness}) então propaga: $x(t + T) =
x(t)$ para todo $t$: periódica.

\textbf{5.} No ramo em que $y > 0$: $y = \sqrt{2(\cos x
- \cos a)}$ e $\dd t = \frac{\dd x}{y}$; integrar $x$
de $-a$ a $a$ dá o meio período, e a simetria $x
\mapsto -x$ divide a integral pela metade de novo:
\[
T(a) = 4\int_0^a\frac{\dd x}{\sqrt{2(\cos x - \cos a)}} .
\]
Convergência em $x = a^-$: $\cos x - \cos a = \sin(a)(a - x) +
O((a-x)^2)$ com $\sin a > 0$: o integrando se comporta como
$\bigl(2\sin a\,(a - x)\bigr)^{-1/2}$, \omterm{def:b3:lebesgue:l1}{integrável}.

\textbf{6.} Com $\sin\frac x2 = k\sin\varphi$, $k =
\sin\frac a2$: $\cos x - \cos a = 2(k^2 - \sin^2\frac x2)
= 2k^2\cos^2\varphi$ e $\frac12\cos\frac x2\,\dd x =
k\cos\varphi\,\dd\varphi$, de modo que
\[
T(a) = 4\int_0^{\pi/2}
\frac{\dd\varphi}{\sqrt{1 - k^2\sin^2\varphi}} .
\]
Quando $a \to 0^+$, $k \to 0$: para $k \leq k_0 < 1$ o
integrando é dominado por $(1 - k_0^2\sin^2\varphi)^{-1/2}$,
\omterm{def:b3:topology:continuity}{contínua} em $\intcc0{\pi/2}$: o TCD dá $T \to
4\cdot\frac\pi2 = 2\pi$. Expandindo $(1 - u)^{-1/2} = 1 +
\frac u2 + \frac{3u^2}8 + \cdots$ com $u =
k^2\sin^2\varphi$ (convergência normal para $k < 1$) e
$\int_0^{\pi/2}\sin^2 = \frac\pi4$:
\[
T(a) = 2\pi\Bigl(1 + \frac{k^2}{4} + O(k^4)\Bigr) :
\]
o isocronismo vale apenas em primeira ordem; o período cresce com
a amplitude.

\textbf{7.} Quando $k \uparrow 1$, os integrandos crescem até $(1
- \sin^2\varphi)^{-1/2} = \frac1{\cos\varphi}$, cuja
integral diverge: por convergência monótona, $T(a) \to
+\infty$ quando $a \to \pi^-$.

\textbf{8.} No ramo $y = 2\cos\frac x2$ ($\abs x <
\pi$): $\frac{\dd x}{2\cos(x/2)} = \dd t$; com $u = \frac
x2$, $\int\frac{\dd u}{\cos u} =
\ln\tan\bigl(\frac u2 + \frac\pi4\bigr)$, de modo que $t =
\ln\tan\bigl(\frac x4 + \frac\pi4\bigr)$, isto é,
\[
x(t) = 4\arctan(\eu^t) - \pi,
\qquad y(t) = x'(t) = \frac{4\eu^t}{1 + \eu^{2t}} =
\frac{2}{\cosh t} .
\]
Verificação pela energia: com $\theta = \arctan\eu^t$,
$\sin2\theta = \frac1{\cosh t}$, de modo que $\cos x = -\cos4\theta =
-1 + \frac{2}{\cosh^2t}$ e $E = \frac{y^2}2 - \cos x =
\frac2{\cosh^2t} + 1 - \frac2{\cosh^2t} = 1$: a trajetória
está na separatriz, e derivar $y^2 = 2(1 +
\cos x)$ onde $y > 0$ reproduz $y' = -\sin x$. Limites: $x
\to \pm\pi$ e $y \to 0$ quando $t \to \pm\infty$.

\textbf{9.} A \omterm{def:b3:groups:action}{órbita} tende à sela $(\pi, 0)$ para a frente
e a $(-\pi, 0)$ para trás, mas nunca chega: se ela atingisse
$(\pi, 0)$ num tempo finito $t^*$, duas \omterm{thm:b3:ode:maximal}{soluções maximais}
distintas --- a solução separatriz e a solução constante
na sela --- passariam pelo mesmo ponto
$(t^*, (\pi, 0))$, contradizendo
o~\cref{cor:b3:ode:uniqueness}. As selas só são atingidas
assintoticamente.

\textbf{10.} Para $E_0 > 1$: $y^2 = 2(E_0 + \cos x) \geq
2(E_0 - 1) > 0$: $y$ mantém seu sinal, e $\abs{x'} = \abs y
\geq \sqrt{2(E_0 - 1)}$: $x$ é estritamente monótona, global
(questão 1), com $x(t) \to \pm\infty$. Como $y(t) =
\pm\sqrt{2(E_0 + \cos x(t))}$ e $\cos$ é $2\pi$-periódico,
$y$ retorna a seu valor cada vez que $x$ avança $2\pi$;
o tempo necessário é
\[
\tau(E_0) = \int_{x_0}^{x_0 + 2\pi}\frac{\dd
x}{\sqrt{2(E_0 + \cos x)}} =
\int_{-\pi}^{\pi}\frac{\dd x}{\sqrt{2(E_0 + \cos x)}}
\]
(substituição; periodicidade): $y$ é $\tau$-periódica ---
o pêndulo gira com taxa de rotação assintoticamente constante
$2\pi/\tau \approx \sqrt{2E_0}$ para energias grandes.

\textbf{11.} O método, em ordem: a \emph{energia} ($\dot E =
0$) reduz o \omterm{ex:b3:ode:planeclassification}{fluxo} bidimensional a curvas de nível
unidimensionais; a \emph{limitação} de $y$ em cada nível, mais a
fuga de \omterm{def:b3:topology:compact}{compactos}, dá a existência global; a
\emph{\omterm{def:b3:topology:compact}{compacidade}} dos níveis fechados dá cotas de velocidade
e, portanto, a periodicidade; a \emph{unicidade} converte o
primeiro retorno em periodicidade exata, proíbe a chegada em tempo
finito às selas e separa os tipos de \omterm{def:b3:groups:action}{órbita};
\emph{linearização e Lyapunov} classificam os equilíbrios;
a \emph{integral do período} é analisada com os teoremas de convergência
do~\cref{ch:b3:lebesgue}. Em ponto algum possuímos --- ou
precisamos de --- uma solução geral em forma fechada: a
teoria qualitativa extraiu da própria equação todas as
características do movimento.

\textbf{12.} Por partes: $W_n = \int_0^{\pi/2}\sin^{2n-1}\varphi
\cdot\sin\varphi\,\dd\varphi = (2n-1)\int_0^{\pi/2}
\sin^{2n-2}\varphi\cos^2\varphi\,\dd\varphi = (2n-1)(W_{n-1}
- W_n)$, de modo que $W_n = \frac{2n-1}{2n}W_{n-1}$; com $W_0 =
\frac\pi2$ e $\prod_{j=1}^n\frac{2j-1}{2j} =
\frac{(2n)!}{4^n(n!)^2}$ (separe $(2n)! = 2^nn!\prod(2j-1)$),
a indução dá o valor em destaque. Série binomial: $(1 -
u)^{-1/2} = \sum_{n\geq0}c_nu^n$ com $c_n =
\frac{(2n)!}{4^n(n!)^2} \in \intoc01$, raio $1$. Para $k
\leq k_0 < 1$ e $u = k^2\sin^2\varphi$, a série
$\sum c_nk^{2n}\sin^{2n}\varphi$ converge normalmente em
$\varphi$ ($c_nk^{2n} \leq k_0^{2n}$), de modo que a integração
termo a termo na fórmula da questão 6 é legítima:
\[
T(a) = 4\sum_{n\geq0}c_nk^{2n}W_n
= 4\cdot\frac\pi2\sum_{n\geq0}c_n^2\,k^{2n}
= 2\pi\sum_{n\geq0}
\Bigl(\frac{(2n)!}{4^n(n!)^2}\Bigr)^{\!2}k^{2n} .
\]
Com $c_0 = 1$, $c_1 = \frac12$, $c_2 = \frac38$: $T(a) =
2\pi\bigl(1 + \frac{k^2}4 + \frac{9k^4}{64} + O(k^6)\bigr)$,
sendo o resto uniforme para $k \leq k_0$ (cauda dominada por uma
série geométrica).

\textbf{13.} $k = \sin\frac a2 = \frac a2 - \frac{a^3}{48} +
O(a^5)$, de modo que
\[
k^2 = \frac{a^2}4 - \frac{a^4}{48} + O(a^6),
\qquad
k^4 = \frac{a^4}{16} + O(a^6) ,
\]
e
\[
\frac{T(a)}{2\pi} = 1 + \frac14\Bigl(\frac{a^2}4 -
\frac{a^4}{48}\Bigr) + \frac9{64}\cdot\frac{a^4}{16} +
O(a^6)
= 1 + \frac{a^2}{16} + \frac{11\,a^4}{3072} + O(a^6) ,
\]
pois $-\frac1{192} + \frac9{1024} = \frac{-16 + 27}{3072}
= \frac{11}{3072}$.

\textbf{14.} Na forma elíptica da questão 6, $a \mapsto
k = \sin\frac a2$ é uma bijeção \omterm{def:b3:topology:continuity}{contínua} estritamente crescente
de $\intoo0\pi$ sobre $\intoo01$ e, para cada
$\varphi \in \intoc0{\pi/2}$, o integrando $(1 -
k^2\sin^2\varphi)^{-1/2}$ é estritamente crescente em $k$: $T$
é estritamente crescente. \omterm{def:b3:topology:continuity}{Continuidade}: em $k \leq k_0 < 1$, o
integrando é dominado pela \omterm{def:b3:topology:continuity}{contínua} $(1 -
k_0^2\sin^2\varphi)^{-1/2}$, de modo que a convergência dominada
se aplica ao longo de $k' \to k$. Com os limites $T \to 2\pi$ quando $a
\to 0^+$ (questão 6) e $T \to +\infty$ quando $a \to \pi^-$
(questão 7), a monotonicidade estrita e o teorema do valor
intermediário fazem de $T$ uma bijeção de $\intoo0\pi$ sobre
$\intoo{2\pi}{+\infty}$.

\textbf{15.} Um relógio conta oscilações; regulado em amplitude
nula, ele contabiliza o período harmônico $2\pi$ por oscilação (na
unidade de tempo do pêndulo). Operando na amplitude $a$, o período
verdadeiro é $T(a) = 2\pi\bigl(1 + \frac{a^2}{16} +
O(a^4)\bigr)$: o relógio contabiliza $2\pi$ enquanto $T(a)$ de fato
transcorre, de modo que ele atrasa pela fração
\[
\frac{T(a) - 2\pi}{T(a)} = \frac{a^2}{16} + O(a^4) .
\]
Para $a = 0.2$ rad (cerca de $11.5$ graus): $\frac{a^2}{16} =
\frac{0.04}{16} = \frac1{400}$, e um dia tem $86400$ s: o
relógio perde $86400/400 = 216$ segundos --- uns três minutos e
meio --- por dia. Daí os dois remédios históricos:
impor uma amplitude minúscula e constante (o escapamento), ou dobrar
o vínculo de modo que o período seja exatamente independente da
amplitude (as faces cicloidais de Huygens, 1657).

\textbf{16.} Em $\tau(E_0) = \int_{-\pi}^{\pi}\frac{\dd
x}{\sqrt{2(E_0 + \cos x)}}$, o integrando é, para cada
$x$ fixo, estritamente decrescente em $E_0$: $\tau$ é estritamente
decrescente. Quando $E_0 \downarrow 1$, os integrandos crescem
pontualmente até $\bigl(2(1 + \cos x)\bigr)^{-1/2} =
\frac1{2\abs{\cos\frac x2}}$, cuja integral em
$\intoo{-\pi}\pi$ diverge ($\cos\frac x2$ se anula em primeira
ordem em $\pm\pi$): a convergência monótona dá $\tau(E_0)
\to +\infty$. Quando $E_0 \to +\infty$:
\[
\sqrt{2E_0}\,\tau(E_0) = \int_{-\pi}^{\pi}
\frac{\dd x}{\sqrt{1 + \cos x/E_0}} \longrightarrow 2\pi
\]
por convergência dominada (para $E_0 \geq 2$ o integrando é
no máximo $\sqrt2$). Logo $\tau \approx 2\pi/\sqrt{2E_0}$: uma
volta leva o tempo da rotação livre à velocidade $\sqrt{2E_0}$,
reduzido o potencial a uma ondulação --- batendo com a taxa de rotação
da questão 10.

\textbf{17.} Os eixos carregam as soluções explícitas $t
\mapsto (x_0\eu^{\alpha t}, 0)$ e $t \mapsto (0,
y_0\eu^{-\gamma t})$, junto com o equilíbrio
$(0, 0)$: eles são uniões de \omterm{def:b3:groups:action}{órbitas}. O campo é $\mathcal
C^1$, logo localmente lipschitziano; uma solução que comece em $Q$
e tocasse um eixo passaria por um ponto de uma
dessas \omterm{def:b3:groups:action}{órbitas} e, pelo~\cref{cor:b3:ode:uniqueness}, coincidiria
com ela --- impossível, uma vivendo no eixo e a outra
não. Logo $Q$ é invariante nas duas direções do tempo.
Equilíbrios em $Q$: $x > 0$ força $\alpha - \beta y = 0$ e
$y > 0$ força $\delta x - \gamma = 0$: o único ponto
$(x_*, y_*) = \bigl(\frac\gamma\delta,
\frac\alpha\beta\bigr)$.

\textbf{18.} Ao longo de uma solução,
\[
\dot H = \Bigl(\delta - \frac\gamma x\Bigr)x' +
\Bigl(\beta - \frac\alpha y\Bigr)y'
= (\delta x - \gamma)(\alpha - \beta y) +
(\beta y - \alpha)(\delta x - \gamma) = 0 .
\]
$f(x) = \delta x - \gamma\ln x$ tem $f'' = \gamma/x^2 > 0$, com
$f'$ anulando-se apenas em $x_*$, e $f \to +\infty$ tanto em
$0^+$ quanto em $+\infty$: estritamente convexa e própria, de mínimo
$f(x_*)$; do mesmo modo, $g(y) = \beta y - \alpha\ln y$, de mínimo
$g(y_*)$. Logo $H \geq h_*$, com igualdade apenas em $(x_*,
y_*)$, e cada subnível $\{H \leq h\}\cap Q$ é \omterm{def:b3:topology:compact}{compacto}:
$f(x) \leq h - g(y_*)$ confina $x$ a um intervalo \omterm{def:b3:topology:compact}{compacto} de
$\intoo0{+\infty}$ por ser própria, do mesmo modo $y$, e o conjunto
é fechado em $\R^2$, pois $H \to +\infty$ no bordo de
$Q$. Uma \omterm{thm:b3:ode:maximal}{solução maximal} permanece em seu conjunto de nível \omterm{def:b3:topology:compact}{compacto}, de modo que
ela não pode deixar todo \omterm{def:b3:topology:compact}{compacto} em tempo finito:
o~\cref{thm:b3:ode:maximal} a torna global.

\textbf{19.} Fixe $h > h_*$ e ponha $c = h - g(y_*) >
f(x_*)$. Como $f$ decresce estritamente de $+\infty$ a
$f(x_*)$ em $\intoc0{x_*}$ e cresce estritamente de volta até
$+\infty$ em $\intco{x_*}{+\infty}$, a equação $f(x) = c$
tem exatamente duas raízes $x_- < x_* < x_+$, e $\{f \leq c\} =
\intcc{x_-}{x_+}$. Para $x \in \intoo{x_-}{x_+}$: $g(y) = h -
f(x) > g(y_*)$ tem exatamente duas raízes $y_-(x) < y_* <
y_+(x)$, \omterm{def:b3:topology:continuity}{contínuas} em $x$ (inversas das restrições
\omterm{def:b3:topology:continuity}{contínuas} estritamente monótonas de $g$ de cada lado de
$y_*$), com $y_\pm(x) \to y_*$ quando $x \to x_\pm$; em $x =
x_\pm$ a única solução é $y = y_*$. Logo $\{H = h\}\cap
Q$ é a união dos gráficos de $y_+$ e $y_-$ sobre
$\intcc{x_-}{x_+}$, colados em $(x_\pm, y_*)$: uma curva fechada
em torno de $(x_*, y_*)$ --- o análogo dos ovais do
pêndulo.

\textbf{20.} Seja $C_h = \{H = h\}\cap Q$ com $h > h_*$: o
único equilíbrio de $Q$ está fora de $C_h$, de modo que o campo nunca
se anula sobre ela. Sinais: $x' = \beta x(y_* - y)$, $y' =
\delta y(x - x_*)$: o movimento vai para a direita abaixo da reta $y
= y_*$, para cima à direita de $x = x_*$, para a esquerda acima e para baixo
à esquerda --- circulação anti-horária. Siga a
solução a partir de um ponto $(x_0, y_-(x_0))$ do ramo inferior aberto,
em que $x' > 0$: o tempo para atingir o canto direito
$B = (x_+, y_*)$ é
\[
\int_{x_0}^{x_+}\frac{\dd x}{\beta x\,\bigl(y_* -
y_-(x)\bigr)} .
\]
Perto de $x_+$, escolha $x' \in \intoo{x_*}{x_+}$; para $x \in
\intcc{x'}{x_+}$, $f(x_+) - f(x) \geq f'(x')\,(x_+ - x)$
($f'$ é crescente e positiva além de $x_*$), ao passo que a
relação de nível e a desigualdade de Taylor dão $g(y_-(x)) -
g(y_*) \leq \frac12\,\bigl(\max g''\bigr)\,(y_* -
y_-(x))^2$ na faixa compacta de $y$ de $C_h$: logo $y_* -
y_-(x) \geq c\,\sqrt{x_+ - x}$ com $c > 0$, e o
integrando é $O\bigl((x_+ - x)^{-1/2}\bigr)$: \omterm{def:b3:lebesgue:l1}{integrável} ---
a convergência da questão 5, transposta. No resto do
ramo, o integrando é \omterm{def:b3:topology:continuity}{contínuo}. Logo $B$ é atingido em
tempo finito; ali $y' = \delta y_*(x_+ - x_*) > 0$, a
\omterm{def:b3:groups:action}{órbita} entra na região $x > x_*$, $y > y_*$, sobe até o
canto superior $(x_*, y_+)$ pela estimativa simétrica (papéis de
$f$ e $g$ trocados), e assim por diante ao longo dos quatro arcos:
após um tempo finito $T > 0$ a solução retorna a seu
ponto de partida. Pelo~\cref{cor:b3:ode:uniqueness}, ela é
$T$-periódica --- o argumento da questão 4, literalmente.

\textbf{21.} Numa \omterm{def:b3:groups:action}{órbita} $T$-periódica em $Q$, $t \mapsto \ln
x(t)$ é $\mathcal C^1$ e $T$-periódica, de modo que
\[
0 = \int_0^T(\ln x)'\,\dd t
= \int_0^T(\alpha - \beta y)\,\dd t
= \alpha T - \beta\int_0^Ty\,\dd t ,
\]
dando $\frac1T\int_0^Ty = \frac\alpha\beta$; do mesmo modo, $0 =
\int_0^T(\ln y)' = \delta\int_0^Tx\,\dd t - \gamma T$ dá
$\frac1T\int_0^Tx = \frac\gamma\delta$. As médias temporais
são os valores de equilíbrio, para toda \omterm{def:b3:groups:action}{órbita}, independentemente da
amplitude --- uma lei de conservação que ninguém colocou à mão.

\textbf{22.} Com colheita, o sistema é de novo da forma de
Lotka--Volterra, com parâmetros $\alpha -
\varepsilon$, $\beta$, $\gamma + \varepsilon$, $\delta$ (o
equilíbrio \omterm{def:b3:topology:interior}{interior} persiste, pois $\varepsilon <
\alpha$). A questão 21 aplicada ao novo sistema:
\[
\bar x = \frac{\gamma + \varepsilon}{\delta}
\ \ (\text{média de presas sobe}),
\qquad
\bar y = \frac{\alpha - \varepsilon}{\beta}
\ \ (\text{média de predadores cai}) :
\]
a colheita indiscriminada desloca o equilíbrio em favor da
presa. Os dados de D'Ancona leem isso ao contrário: a guerra cortou
a pesca, $\varepsilon$ caiu, de modo que a média de predadores
$(\alpha - \varepsilon)/\beta$ subiu e a média de presas
caiu --- uma fração maior de tubarões na captura, exatamente o que
os mercados de peixe do Adriático registraram. Esse é o princípio de
Volterra, o mesmo mecanismo por trás dos paradoxos dos pesticidas:
abater os dois níveis tróficos beneficia o nível que é comido.

\textbf{23.} A receita, nas duas vezes: (i) uma integral primeira
--- $E$ para o pêndulo, $H$ aqui, obtida separando
$\dd y/\dd x$ --- colapsa o plano em curvas; (ii) a
propriedade de ser própria e a \omterm{def:b3:topology:compact}{compacidade} dos conjuntos de nível dão a existência
global pelo~\cref{thm:b3:ode:maximal}; (iii) a
geometria dos níveis --- ovais, vindos da forma de $\cos$
lá e da convexidade estrita de $f$ e $g$ aqui --- é
lida na integral, não no \omterm{ex:b3:ode:planeclassification}{fluxo}; (iv) um campo não nulo
num oval \omterm{def:b3:topology:compact}{compacto}, mais singularidades \omterm{def:b3:lebesgue:l1}{integráveis} nos cantos,
força um tempo de retorno finito; (v) a unicidade
(\cref{cor:b3:ode:uniqueness}) converte retorno em
periodicidade e proíbe a chegada em tempo finito aos equilíbrios;
(vi) os dividendos --- expansões do período, leis de médias
--- vêm dos teoremas de convergência aplicados às
integrais resultantes. Nem $x'' = -\sin x$ nem
Lotka--Volterra admite solução elementar em forma fechada
(integrais elípticas num caso, curvas de nível transcendentes
no outro) e, em ponto algum, foi preciso uma: a própria
equação, interrogada qualitativamente, entregou o movimento
inteiro.

\textbf{24.} Num nível de energia $E_0 > 1$, $y^2 = 2(E_0 +
\cos x) > 0$: extremos de $y^2$ em $\cos x = \pm1$, dando
as cotas enunciadas, atingidas em $x \equiv 0, \pi$. Média
temporal num período $\tau = \tau(E_0)$: $x$ avança
exatamente $2\pi$, de modo que
\[
\frac1\tau\int_0^\tau y\,\dd t = \frac{x(\tau) -
x(0)}{\tau} = \frac{2\pi}\tau .
\]
A razão entre as velocidades extremas é $\sqrt{\frac{E_0 + 1}{E_0 -
1}} = 1 + O(E_0^{-1}) \to 1$: em alta energia, a ondulação
$\pm1$ do potencial é desprezível diante de $E_0$, e o pêndulo
gira quase uniformemente --- a tábua de lavar se aplaina.

\textbf{25.} Em $E_0 \geq 1 + \delta$, o integrando de
$\tau(E_0) = \int_{-\pi}^{\pi}\frac{\dd x}{\sqrt{2(E_0 +
\cos x)}}$ é dominado por $(2\delta)^{-1/2}$, e sua
derivada em $E_0$,
\[
\partial_{E_0}\frac1{\sqrt{2(E_0 + \cos x)}}
= -\frac{1}{\bigl(2(E_0 + \cos x)\bigr)^{3/2}}
\]
por $(2\delta)^{-3/2}$, ambos \omterm{def:b3:lebesgue:l1}{integráveis} no \omterm{def:b3:topology:compact}{compacto}
$\intcc{-\pi}\pi$: a derivação sob o sinal de integral
(\cref{thm:b3:lebesgue:paramdiff}) se aplica e dá
$\tau'(E_0) < 0$ (o integrando é estritamente negativo):
estritamente decrescente, $\mathcal C^1$. Limites: quando $E_0 \to
\infty$, $\tau \leq \frac{2\pi}{\sqrt{2(E_0 - 1)}} \to 0$;
quando $E_0 \to 1^+$: escreva $E_0 + \cos x = (E_0 - 1) +
2\cos^2\frac x2$; o integrando cresce quando $E_0$
decresce, de modo que, por convergência monótona,
\[
\tau(E_0) \;\nearrow\;
\int_{-\pi}^{\pi}\frac{\dd x}{2\,\abs{\cos\frac x2}}
= +\infty,
\]
divergindo a integral limite em $x = \pm\pi$ (ali
$\abs{\cos\frac x2} \sim \frac{\abs{x \mp \pi}}2$, um
$\frac1{\abs\cdot}$ não \omterm{def:b3:lebesgue:l1}{integrável}): o período explode
ao se aproximar da separatriz, batendo com o $T(a) \to
\infty$ da Parte II pelo lado das librações. O eixo
da energia se lê: repouso em $E_0 = -1$; librações, $2\pi \nearrow
\infty$, em $\intoo{-1}1$; a separatriz infinitamente lenta em
$E_0 = 1$; rotações, $\infty \searrow 0$, além. Uma
integral, a vida inteira do pêndulo.
\end{solution}
